% Created 2023-01-28 Sat 19:31
% Intended LaTeX compiler: pdflatex
\documentclass[11pt]{article}
\usepackage[utf8]{inputenc}
\usepackage[T1]{fontenc}
\usepackage{graphicx}
\usepackage{longtable}
\usepackage{wrapfig}
\usepackage{rotating}
\usepackage[normalem]{ulem}
\usepackage{amsmath}
\usepackage{amssymb}
\usepackage{capt-of}
\usepackage{hyperref}
\author{mahmooz}
\date{\today}
\title{random variable}
\hypersetup{
 pdfauthor={mahmooz},
 pdftitle={random variable},
 pdfkeywords={},
 pdfsubject={},
 pdfcreator={Emacs 30.0.50 (Org mode 9.6)}, 
 pdflang={English}}
\begin{document}

\maketitle
\tableofcontents

\begin{definition}
a \textbf{random variable} X is a function \(X(\Omega)\) from the \href{probability.org}{sample space} to the real line (\(\mathbb{R}\)), such that every event in the sample space corresponds to a real number. \(P_i\) is the probability of receiving the value \(X_i\) and \(P(X)\) is the probability function\\[0pt]
there are two types of random variables\\[0pt]
\begin{itemize}
\item \href{20230121192455-discrete_random_variable.org}{discrete random variable}\\[0pt]
\item \href{20230128192014-continuous_random_variable.org}{continuous random variable}\\[0pt]
\end{itemize}
\begin{my_example}
we roll a symmetric coin 3 times\\[0pt]
\[
\Omega = \{(h,h,h),(h,h,t),(h,t,h),(t,h,h),(t,h,t),(h,t,t),(t,t,h),(t,t,t)\}
\]\\[0pt]
we define \(X\) as the number of times that "heads" was received\\[0pt]
\(X\) is a random variable with 4 possible values, \(X(\Omega)=0,1,2,3\), \(P(X=0)=\frac{1}{8},P(X=1)=\frac{3}{8},P(X=2)=\frac{3}{8},P(X=3)=\frac{1}{8}\)\\[0pt]
\end{my_example}
\begin{lemma}
let \(f_i\) be the \href{20230113185822-frequency.org}{frequency} of the value \(x_i\), we get \(\sum_{i}^{k} f_i = n\), so \(\overline{x} = \frac{\sum_{i=1}^{k} x_if_i}{n}\)\\[0pt]
let \(p_i=\frac{f_i}{n}\) be the \href{probability.org}{probability} to get the value \(x_i\), so \(\sum_{i}^{k} p_i=1\), we get \(\overline{x}=\frac{\sum_{1}^{k} x_if_i}{n}=\frac{1}{n}\sum_{i=1}^{k} x_if_i=\sum_{i=1}^{k} x_i \cdot \frac{f_i}{n}=\sum_{i=1}^{k} x_ip_i\)\\[0pt]
\[
\overline{x} = \sum_{i=1}^{k} x_ip_i
\]\\[0pt]
\end{lemma}
\begin{lemma}
let \(x_1,x_2,\dots,x_k\) be arbitrary values and \(f_1,f_2,\dots,f_k\) their frequencies\\[0pt]
the \href{20230113164748-variance.org}{variance}s of \(x_1,x_2,\dots,x_k\) is the measure of their scattering around the \href{20220813152209-average.org}{average}\\[0pt]
\[
S^2 = \frac{1}{n} \sum_{i=1}^{k} (x_i-\overline{x})^2 f_i
\]\\[0pt]
\end{lemma}
\begin{lemma}
we build off from the previous lemmas:\\[0pt]
\begin{align*}
S^2 &= \frac{1}{n} \sum_{i=1}^{k} (x_i-\overline{x})^2 f_i\\
&= \sum_{i=1}^{k} (x_i^2-2x_i\overline{x}+\overline{x}^2) \cdot \frac{f_i}{n}\\
&= \sum_{i=1}^{k} (x_i^2-2x_i\overline{x}+\overline{x}^2) \cdot p_i\\
&= \sum_{i=1}^{k} x_i^2p_i - 2\sum_{i=1}^{k} x_ip_i\overline{x} + \sum_{i=1}^{k} \overline{x}^2p_i\\
&= \sum_{i=1}^{k} x_i^2p_i - 2\overline{x}\sum_{i=1}^{k} x_ip_i + \overline{x}^2\sum_{i=1}^{k} p_i\\
&= \sum_{i=1}^{k} x_i^2p_i - 2\overline{x}^2 + \overline{x}^2\\
&= \sum_{i=1}^{k} x_i^2p_i - \overline{x}^2
\end{align*}
so\\[0pt]
\[
S^2 = \sum_{i=1}^{k} x_i^2p_i - \overline{x}^2
\]\\[0pt]
\end{lemma}
\begin{lemma}
\(E(x)\) is the \href{20230121184721-expected_value.org}{expected value} of a random variable \(X\), which is the average of it, so we get \(\overline{x} = \sum_{i=1}^{k} x_ip_i\) so \(E(x) = \sum_{i=1}^{k} x_ip_i\)\\[0pt]
\end{lemma}
\begin{my_example}
on his way to work a person goes through 3 traffic lights, the probability of the first being green is 0.6, the second 0.5, the third 0.4\\[0pt]
let \(x\) be the number of traffic lights they pass without waiting (green)\\[0pt]
\begin{subquestion}
find the probability function of the random variable \(x\)\\[0pt]
\begin{answer}
\(x\) can equal one of \(0,1,2,3\)\\[0pt]
\(p(x=0)=0.4 \cdot 0.5 \cdot 0.6 = 0.12\) (no green)\\[0pt]
\(p(x=1)=0.6 \cdot 0.6 \cdot 0.5 \cdot 0.6 + 0.4 \cdot 0.5 \cdot 0.6 + 0.4 \cdot 0.5 \cdot 0.4 = 0.38\) (1 green)\\[0pt]
\(p(x=2)=0.6 \cdot 0.5 \cdot 0.6 + 0.4 \cdot 0.5 \cdot 0.4 + 0.6 \cdot 0.5 \cdot 0.4 = 0.38\) (2 green)\\[0pt]
\(p(x=3)=0.6 \cdot 0.5 \cdot 0.4 = 0.12\) (3 green)\\[0pt]
so \(p(x) = \begin{cases} 0.12 \quad x=0\\ 0.38 \quad x=1\\ 0.38 \quad x=2\\ 0.12 \quad x=3 \end{cases}\)\\[0pt]
\end{answer}
\end{subquestion}
\begin{subquestion}
find the expected value and the variance of \(x\)\\[0pt]
\begin{answer}
\(\mathrm{E}(x)=0.12 \cdot 0 + 0.38 \cdot 1 + 0.38 \cdot 2 + 0.12 \cdot 3 = 1.5\)\\[0pt]
\(\mathrm{E}(x)=0.12 \cdot 0^2 + 0.38 \cdot 1^2 + 0.38 \cdot 2^2 + 0.12 \cdot 3^2 = 2.98\)\\[0pt]
\(\mathrm{Var}(x)=\mathrm{E}(x^2)-\mathrm{E}(x)^2=2.98-(1.5)^2=2.98-2.25=0.73\)\\[0pt]
\end{answer}
\end{subquestion}
\end{my_example}
\end{definition}
\end{document}